\section{自由群的子群}
\label{ch14-sec85}

现在我们来证明主要定理, 即关于自由群 $F$ 的子群 $H$ 还是自由群等相关结论. 证明中用到的方法可圈可点, 当 $F$ 的自由生成元组的基数已知时, 这个方法使我们对于 $H$ 的自由生成元组的基数有所了解.

\begin{theorem}[自由群的子群是自由群]
\label{theorem:ch14-subgroup-free}
如果 $H$ 是自由群 $F$ 的子群, 则 $H$ 是一个自由群.
\end{theorem}

\begin{proof}
设 $\{\alpha\mid\alpha\in J\}$ 是 $F$ 的一个自由生成元组. 设 $X$ 是一些圆周 $S_{\alpha}$ ($\alpha\in J$) 的束, $x_{0}$ 是这些圆周的公共点. 我们按以下方式将 $X$ 构造成一个线性图: 将每一个圆周 $S_{\alpha}$ 分为三条弧, 其中两条以 $x_{0}$ 为端点. 对于每一个 $\alpha$ 指定一条生成 $\pi_{1}(S_{\alpha},\ x_{0})$ 的回路, 这诱导出 $F$ 与 $\pi_{1}(X,\ x_{0})$ 之间的一个同构. 因此, 我们假设 $F$ 等于群 $\pi_{1}(X,\ x_{0})$.

空间 $X$ 是道路连通、局部道路连通且半局部单连通的. 应用 \cref{theorem:ch13-covering-existence} 可见, 存在 $X$ 的一个道路连通的覆叠空间 $p\colon E\to X$, 使得对于 $p^{-1}(x_{0})$ 中的某一个点 $e_{0}$,
\[
p_{*}(\pi_{1}(E,e_{0}))=H.
\]
由于 $p_{*}$ 是一个单同态, 所以 $\pi_{1}(E,\ e_{0})$ 同构于 $H$.

根据 \cref{theorem:ch14-covering-of-graph}, 空间 $E$ 是一个线性图. 于是 \cref{theorem:ch14-graph-fundamental-free} 蕴涵着其基本群是一个自由群.
\end{proof}

\begin{definition}[Euler 数]
\label{definition:ch14-euler-number}
如果 $X$ 是一个有限线性图, 定义 $X$ 的\textbf{Euler 数}(Euler number) 为 $X$ 的顶点个数减去边的个数. 用希腊字母 $\chi$ 来表示, 如 $\chi(X)$.
\end{definition}

\begin{lemma}[有限连通图的秩公式]
\label{lemma:ch14-rank-formula}
如果 $X$ 是有限的、连通的线性图, 则 $X$ 的基本群的自由生成元组的基数为 $1-\chi(X)$.
\end{lemma}

\begin{proof}
\textbf{第一步.} 我们先证明对于任何有限的树 $T$ 有 $\chi(T)=1$. 对 $T$ 中所含边的数目 $n$ 施行归纳法. 当 $n=1$ 时, $T$ 有一条边和两个顶点, 所以 $\chi(T)=1$. 如果 $n>1$, 则 $T=T_{0}\cup A$, 其中 $T_{0}$ 是所含边的数目为 $n-1$ 的一棵树, $A$ 是与 $T_{0}$ 仅交于一个顶点的一条边. 根据归纳假设, 我们有 $\chi(T_{0})=1$. 图 $T$ 比 $T_{0}$ 多一个顶点和一条边, 所以 $\chi(T)=1$.

\textbf{第二步.} 现在来证明定理. 给定 $X$, 设 $T$ 为 $X$ 中的一棵极大树. 如果 $X=T$, 则已经完成了证明. 如若不然, 设 $A_{1}$, $\cdots$, $A_{n}$ 是 $X$ 的那些不在 $T$ 中的边. 这时 $X$ 的基本群有一个 $n$ 个元素的自由生成元组. 另一方面, $X$ 与 $T$ 有相同的顶点数, 并且 $X$ 比 $T$ 多了 $n$ 条边. 因此
\[
\chi(X)=\chi(T)-n=1-n,
\]
所以 $n=1-\chi(X)$.
\end{proof}

\begin{definition}[指数]
\label{definition:ch14-index}
设 $H$ 是群 $G$ 的一个子群. 如果 $H$ 在 $G$ 中的右陪集族 $G/H$ 是有限的, 则其基数称为 $H$ 在 $G$ 中的\textbf{指数}(index). (当然, $H$ 在 $G$ 中的左陪集具有相同的基数.)
\end{definition}

\begin{theorem}[Schreier 指数公式]
\label{theorem:ch14-schreier}
设 $F$ 是一个具有 $n+1$ 个自由生成元的自由群, $H$ 是 $F$ 的一个子群. 如果 $H$ 在 $F$ 中的指数为 $k$, 则 $H$ 有 $kn+1$ 个自由生成元.
\end{theorem}

\begin{proof}
我们运用 \cref{theorem:ch14-subgroup-free} 中的做法. 假设 $F=\pi_{1}(X,\ x_{0})$, 其中 $X$ 是一个线性图, 其底空间是 $n+1$ 个圆周的束. 给定 $H$, 选取一个道路连通的覆叠空间 $p\colon E\to X$, 使得 $p_{*}(\pi_{1}(E,\ e_{0}))=H$. 此时提升对应
\[
\Phi\colon\pi_{1}(X,x_{0})/H\longrightarrow p^{-1}(x_{0})
\]
是一个一一映射. 因此, $E$ 是 $X$ 的一个 $k$-重覆叠映射.

空间 $E$ 也是一个线性图. 给定 $X$ 中的一条边 $A$, $p^{-1}(A)$ 的道路分支都是 $E$ 的边, 并且这些边中的每一条都被 $p$ 同胚地映到 $A$ 上. 于是 $E$ 的边数是 $X$ 的边数的 $k$ 倍, 顶点也是 $k$ 倍. 这蕴涵着 $\chi(E)=k\chi(X)$. 由于 $X$ 的基本群有 $n+1$ 个自由生成元, 前面一个引理告诉我们 $\chi(X)=-n$. 于是 $E$ 的基本群的自由生成元的个数是
\[
1-\chi(E)=1-k\chi(X)=1+kn,
\]
因为 $E$ 的基本群同构于 $H$.
\end{proof}

注意, 如果 $F$ 是一个具有有限自由生成元组的自由群, 而 $H$ 是 $F$ 的一个子群, 使得 $F/H$ 是无限的, 那么关于 $H$ 的自由生成元组的基数我们得不到什么结论. 它可能是有限的 (例如, 当 $H$ 是平凡子群时), 也可能是无限的 (例如, 当 $H$ 是画在第 81 节例 2 中的那个覆叠空间的基本群时).

\begin{exercise}
\label{exercise:ch14-sec85-01}
证明有限线性图 $X$ 的 Euler 数是 $X$ 的一个拓扑不变量. [提示: 先考虑 $X$ 连通的情形.]
\end{exercise}

\begin{exercise}
\label{exercise:ch14-sec85-02}
设 $F$ 是由两个生自由成元 $\alpha$ 和 $\beta$ 生成的自由群, 并设 $H$ 是由 $\alpha$ 生成的子群. 证明 $H$ 在 $F$ 中的指数是无限的.
\end{exercise}

\begin{exercise}
\label{exercise:ch14-sec85-03}
设 $p\colon\mathbb{R}\to S^{1}$ 是标准的覆叠映射, 考虑覆叠映射 $p\times p\colon\mathbb{R}\times\mathbb{R}\to S^{1}\times S^{1}$. 令 $b_{0}=(1,\ 0)\in S^{1}$, 设 $X=(b_{0}\times S^{1})\cup(S^{1}\times b_{0})$, $E=(p\times p)^{-1}(X)$. 令 $q\colon E\to X$ 是限制 $p\times p$ 得到的覆叠映射. $X$ 的基本群有两个生成元 $\alpha$ 和 $\beta$, 其中 $\alpha$ 由 $b_{0}\times S^{1}$ 中的一条回路代表, $\beta$ 由 $S^{1}\times b_{0}$ 中的一条回路代表. 找出子群 $q_{*}(\pi_{1}(E,\ e_{0}))$ 的一个自由生成元组, 其中 $e_{0}$ 是 $\mathbb{R}^{2}$ 中的原点.
\end{exercise}
